\subsection{Reverse-KL training algorithm}

At prefix $s_t$, OPD asks the student's next-token distribution to approach the teacher's
distribution at the same prefix. Thinking Machines Lab uses reverse KL:
\begin{equation}
  D_{\mathrm{KL}}\!\left(\pi_\theta(\cdot\mid s_t)
  \,\|\,\pi_T(\cdot\mid s_t)\right)
  =\mathbb{E}_{a\sim\pi_\theta}
  \left[\log\pi_\theta(a\mid s_t)-\log\pi_T(a\mid s_t)\right].
  \label{eq:opd-reverse-kl}
\end{equation}
The value is zero when the next-token distributions match. Reverse KL is mode-seeking: it
concentrates the student on behavior the teacher considers likely.

The implementation does not need to sample a new answer from the teacher. It performs four
steps:
\begin{enumerate}
  \item Sample exact token trajectories from a student snapshot $\pi_{\mathrm{old}}$ and
        retain the sampled tokens' log-probabilities.
  \item Run the teacher on the same prefixes and record its log-probability for each
        student-sampled token.
  \item Set $A_t=\log\pi_T(a_t\mid s_t)-\log\pi_{\mathrm{old}}(a_t\mid s_t)$: teacher-like
        tokens receive positive signal and implausible tokens receive negative signal.
  \item Update only the student. For reused rollouts, weight $A_t$ by the RL importance
        ratio $\pi_\theta(a_t\mid s_t)/\pi_{\mathrm{old}}(a_t\mid s_t)$.
\end{enumerate}
Prompt, padding, and tool-response tokens remain conditioning context but are masked from
the policy loss. Unlike a terminal advantage broadcast over a whole answer, $A_t$ changes
at every trained token.
